% =====================================================================
% Fig.1 -- Overall architecture of DPRNet.
% Drawn in the convention used by lightweight-SR papers (SwinIR, CATANet):
% the data path is rendered as feature-map slabs, the deep extractor is a
% stacked "block x L" slab, and sub-modules are expanded through shaded
% insets rather than hairline zoom lines. All equations stay in the text.
% =====================================================================
\begin{figure*}[t]
\centering
% --- feature-map slab: name, cx, cy, w, h, front, top, side ---
\newcommand{\fmslab}[8]{%
  \begin{scope}[shift={(#2,#3)}]
    \pgfmathsetmacro\hw{#4/2}\pgfmathsetmacro\hh{#5/2}%
    \fill[#6] (-\hw,-\hh) rectangle (\hw,\hh);
    \fill[#7] (-\hw,\hh)--(-\hw+2,\hh+2)--(\hw+2,\hh+2)--(\hw,\hh)--cycle;
    \fill[#8] (\hw,-\hh)--(\hw+2,-\hh+2)--(\hw+2,\hh+2)--(\hw,\hh)--cycle;
    \draw[black!65,line width=0.5pt] (-\hw,-\hh) rectangle (\hw,\hh);
    \draw[black!65,line width=0.5pt]
      (-\hw,\hh)--(-\hw+2,\hh+2)--(\hw+2,\hh+2)--(\hw+2,-\hh+2)--(\hw,-\hh);
    \draw[black!65,line width=0.5pt] (\hw,\hh)--(\hw+2,\hh+2);
    \coordinate (#1) at (0,0);
    \coordinate (#1E) at (\hw+2,0);
    \coordinate (#1W) at (-\hw,0);
    \coordinate (#1N) at (1,\hh+1);
    \coordinate (#1S) at (0,-\hh);
  \end{scope}}
\resizebox{0.98\textwidth}{!}{%
\begin{tikzpicture}[
    x=1mm, y=1mm,
    font=\footnotesize,
    >={Latex[length=4pt]},
    box/.style={rounded corners=2pt, draw=black!75, fill=black!7,
                minimum height=8mm, align=center, inner sep=2.5pt,
                font=\scriptsize},
    data/.style={rounded corners=2pt, draw=dpramber!85!black, line width=0.7pt,
                 fill=dpramber!20, minimum height=7mm, align=center,
                 inner sep=2.2pt, font=\scriptsize},
    dprbox/.style={rounded corners=2pt, draw=dprblue, line width=1pt,
                   fill=dprblue!22, minimum height=8mm, align=center,
                   inner sep=2.5pt, font=\scriptsize},
    sum/.style={circle, draw=black!75, fill=white, inner sep=0pt,
                minimum size=4.6mm, font=\scriptsize},
    flow/.style={->, draw=black!82, line width=0.7pt},
    skip/.style={->, draw=black!55, line width=0.6pt, dashed},
    lab/.style={font=\scriptsize, black!75, inner sep=1.5pt},
    note/.style={font=\scriptsize, black!65},
    tag/.style={font=\footnotesize\bfseries, black!80},
    panel/.style={rounded corners=3pt, draw=black!28, fill=white,
                  inner sep=3.5mm}
]

% ===================== (a) Network backbone =====================
\fmslab{img}{5}{0}{8}{13}{dpramber!22}{dpramber!36}{dpramber!50}
\node[box, minimum width=15mm] (conv) at (24,0) {Conv $3{\times}3$};
\fmslab{ff0}{41}{0}{6}{11}{dpramber!18}{dpramber!32}{dpramber!46}
% --- stacked deep-feature blocks ---
\foreach \i in {2,1,0}{%
  \begin{scope}[shift={($(64,0)+(1.8*\i,1.8*\i)$)}]
    \fill[dprblue!16] (-11,-7) rectangle (11,7);
    \fill[dprblue!30] (-11,7)--(-9,9)--(13,9)--(11,7)--cycle;
    \fill[dprblue!40] (11,-7)--(13,-5)--(13,9)--(11,7)--cycle;
    \draw[dprblue!70,line width=0.55pt] (-11,-7) rectangle (11,7);
    \draw[dprblue!70,line width=0.55pt] (-11,7)--(-9,9)--(13,9)--(13,-5)--(11,-7);
    \draw[dprblue!70,line width=0.55pt] (11,7)--(13,9);
  \end{scope}}
\node[align=center, font=\scriptsize, dprblue!85!black] at (64,0)
  {\textbf{DPRNet}\\\textbf{block} $\times L$};
\coordinate (blkW) at (53,0);
\coordinate (blkE) at (77,0);
\fmslab{ffl}{91}{0}{6}{11}{dpramber!18}{dpramber!32}{dpramber!46}
\node[sum] (gsum) at (104,0) {$+$};
\node[box, minimum width=24mm] (up) at (124,0) {Pixel-shuffle\\upsample};
\node[sum] (osum) at (146,0) {$+$};
\fmslab{sr}{160}{0}{8}{13}{dpramber!22}{dpramber!36}{dpramber!50}

% labels under the I/O slabs
\node[lab, below=0.5mm of imgS] {$I_{LR}$};
\node[lab, below=0.5mm of srS] {$I_{SR}$};
% coordinates that let the background panel enclose the 3D slab/stack tops
\coordinate (fitTL) at (0,9);
\coordinate (fitTM) at (68,13);
\coordinate (fitTR) at (164,9);

% main flow
\draw[flow] (imgE) -- (conv.west);
\draw[flow] (conv.east) -- node[lab,above]{$F_0$} (ff0W);
\draw[flow] (ff0E) -- (blkW);
\draw[flow] (blkE) -- node[lab,above]{$F_L$} (fflW);
\draw[flow] (fflE) -- (gsum.west);
\draw[flow] (gsum.east) -- (up.west);
\draw[flow] (up.east) -- (osum.west);
\draw[flow] (osum.east) -- (srW);
% global residual (F_0) and image bypass (I_LR)
\draw[skip] (ff0S) |- (41,-13) -| node[pos=0.25,below,lab]{global residual} (gsum.south);
\draw[skip] (imgS) |- (5,-20) -| node[pos=0.22,below,lab]{bilinear $\uparrow s$} (osum.south);
\node[tag, anchor=west] at (-3,11) {(a)};

% ===================== (b) One DPRNet block =====================
\node[data] (finb) at (10,-44) {$F_{l-1}$};
\node[dprbox, minimum width=15mm] (tab) at (28,-44) {TAB$_l$};
\node[box, minimum width=15mm] (lrsa) at (47,-44) {LRSA$_l$};
\node[box, minimum width=15mm] (fconv) at (65,-44) {Conv $3{\times}3$};
\node[sum] (badd) at (78,-44) {$+$};
\node[data] (flo) at (88,-44) {$F_l$};

\draw[flow] (finb) -- (tab);
\draw[flow] (tab) -- (lrsa);
\draw[flow] (lrsa) -- (fconv);
\draw[flow] (fconv) -- (badd);
\draw[flow] (badd) -- (flo);
\draw[skip] (finb.south) |- (10,-55) -| node[pos=0.35,below,lab]{block residual} (badd.south);
\node[tag, anchor=west] at (-3,-34) {(b)};

% ===================== (c) Inside TAB =====================
\node[data] (xin) at (97,-44) {$X$};
\node[box, minimum width=11mm] (ln) at (108,-44) {LN};
\node[dprbox, minimum width=12mm] (dpr) at (122,-44) {DPR};
\node[box, minimum width=12mm] (iasa) at (137,-44) {IASA};
\node[box, minimum width=13mm] (c11) at (152,-44) {Conv $1{\times}1$};
\node[box, minimum width=14mm] (ffn) at (168,-44) {ConvFFN};
\node[data] (tabout) at (180,-44) {out};

\draw[flow] (xin) -- (ln);
\draw[flow] (ln) -- (dpr);
\draw[flow] (dpr) -- node[lab,above,font=\tiny]{Fig.~\ref{fig:dpr}} (iasa);
\draw[flow] (iasa) -- (c11);
\draw[flow] (c11) -- (ffn);
\draw[flow] (ffn) -- (tabout);
\node[tag, anchor=west] at (90,-34) {(c)};

% ===================== panels + expansion wedges =====================
\begin{scope}[on background layer]
  \node[panel, fit=(img)(conv)(blkW)(blkE)(up)(osum)(sr)(gsum)(fitTL)(fitTM)(fitTR)] (pa) {};
  \node[panel, fit=(finb)(tab)(lrsa)(fconv)(badd)(flo)] (pb) {};
  \node[panel, fit=(xin)(ln)(dpr)(iasa)(c11)(ffn)(tabout)] (pc) {};
  % (a) deep-feature stack  ->  (b) one block
  \fill[dprblue!7] (53,-7) -- (75,-7) -- (pb.north east) -- (pb.north west) -- cycle;
  \draw[dprblue!40, dashed, line width=0.5pt] (53,-7) -- (pb.north west);
  \draw[dprblue!40, dashed, line width=0.5pt] (75,-7) -- (pb.north east);
  % (b) TAB  ->  (c) inside TAB
  \fill[dprblue!7] (tab.south west) -- (tab.south east) -- (pc.north west) -- ($(pc.north west)!0.6!(pc.north east)$) -- cycle;
  \draw[dprblue!40, dashed, line width=0.5pt] (tab.south west) -- (pc.north west);
  \draw[dprblue!40, dashed, line width=0.5pt] (tab.south east) -- ($(pc.north west)!0.6!(pc.north east)$);
\end{scope}

\end{tikzpicture}%
}
\caption{Overall architecture of DPRNet, following the lightweight-SR backbone of
CATANet. \textbf{(a)}~A $3{\times}3$ convolution extracts shallow features $F_0$,
which pass through a stack of $L{=}8$ DPRNet blocks; the deep features are added
back through a global residual and reconstructed by a pixel-shuffle upsampler that
is summed with a bilinearly upsampled base image. \textbf{(b)}~Each block keeps the
TAB--LRSA--fusion structure under a residual connection. \textbf{(c)}~Inside a TAB,
the only modified component is the routing module: DPR replaces center-based
routing while preserving the input/output interface, and is expanded in
Fig.~\ref{fig:dpr}.}
\label{fig:arch}
\end{figure*}
